\section{Chebyshev Polynomials}
\label{sec:Chebyshev polynomials}
The Chebyshev polynomials of the first kind are the eigenfunctions of the singular Sturm--Liouville problem \eqref{eq:Sturm-Liouville} with $p(x) = \sqrt{1-x^2}$, $q = 0$, and $w = (1-x^2)^{-\frac 1 2}$
\begin{equation*}
    -\frac{\diff}{\diff x} \left( \sqrt{1-x^2} \frac{\diff T_k}{\diff x} \right) =  \frac{k^2}{\sqrt{1-x^2}} T_k(x).
\end{equation*}

\begin{remark}
    The reason for choosing $p$ degenerate at $x = \pm 1$ (i.e., $p(\pm 1) = 0$) is justified by \Cref{rem:Sturm-Liouville degenerate}.
    The choice of a weight $w$ that is singular at $x = \pm 1$ can be explained by \Cref{sec:Runge's phenomenon}, although $P_N u$ is not the interpolation.
    It is also common in this framework that most of the error happens near the boundary $x = \pm 1$.
    We expect the $L^\infty$-error to improve if we work with $L^2_w$-error for $w$ singular on the boundary instead of $L^2$-error.
\end{remark}


The solution is precisely
\begin{equation}
    \label{eq:Chebyshev Tn}
    T_k (x) = \cos \left( k \arccos x \right).
\end{equation}
We have the recurrence relation
\begin{equation}
    \label{eq:Chebyshev recursion}
    T_0(x) = 1,
    \qquad
    T_1(x) = x,
    \qquad
    T_{n+1} (x) = 2x T_n(x) - T_{n-1}(x).
\end{equation}

\subsection{Quadrature rules}

\begin{example}[Gauss integration for the Chebyshev weight $w = (1 - x^2)^{-\frac 1 2}$]
    Due to \Cref{rem:Sturm-Liouville degenerate} we can apply \Cref{thm:Gauss integration} with
    \begin{equation*}
        x_j = \cos \frac{(2j+1)\pi}{2N+2} , \qquad w_j = \frac{\pi}{N+1}.
    \end{equation*}
\end{example}

We have that
\begin{equation*}
    \|u - P_N u\|_{L^2_w (-1,1)} \le C N^{-m} \| u \|_{H^m_w (-1,1)},
\end{equation*}
where $H^m_w$ is the weighted Sobolev space. The same holds for $I_N$ at the Gauss or Gauss--Lobatto nodes.

\begin{example}[Gauss--Lobatto integration for the Chebyshev weight]
    \label{ex:Chebyshev Gauss-Lobatto}
    Let $w = (1 - x^2)^{-\frac 1 2}$, i.e., $\alpha = \beta = -\frac 1 2$.
    Applying \Cref{thm:Gauss-Lobatto integration Jacobi weights} we get precisely
    \begin{equation*}
        x_k = \cos \frac{k \pi}{N} \text{ for } k \in \{0, \dots, N \},
    \end{equation*}
    and
    \begin{equation*}
        w_0 = w_N = \frac{\pi}{2N}, \qquad w_k = \frac{\pi}N \text{ for } k \in \{1, \dots, N-1\}.
    \end{equation*}
\end{example}

The Chebyshev polynomials of the second kind are given by
\begin{equation*}
    U_n (x) = \frac{\sin \big( (n+1) \arccos x \big)}{\sin \left( \arccos x \right)}.
\end{equation*}
We have that
\begin{equation}
    \label{eq:Chebyshev derivatives}
    \frac{\diff T_n}{\diff x} = n U_{n-1}, \qquad \frac{\diff U_n}{\diff x} = \frac{(n+1) T_{n+1} - x U_n}{x^2 - 1}, \qquad \frac{\diff^2T_n}{\diff x^2} = n \frac{n T_n - x U_{n-1}}{x^2 - 1}.
\end{equation}

\subsection{Differentiation}

When the $x_j$ are the Gauss--Lobatto points of \namedCref{ex:Chebyshev Gauss-Lobatto}, the coefficients of the differentiation matrices $\mat D_N$ of \Cref{sec:spectral collocation linear algebra} can be found explicitly in \cite[Section 2.4]{CanutoHussainiQuarteroniZang2011}.
